% !TeX root = ../thuthesis-example.tex

\chapter{系统分析设计}

\section{需求分析}

\subsection{功能需求}

本项目面向类即时通信与内容社区系统，核心功能需求如下：

\begin{enumerate}
  \item 用户与关系管理。系统需要支持用户资料、联系人关系、关注关系、群组成员关系等数据管理，为推荐系统提供社交图谱基础。
  \item 内容管理。系统需要保存新闻内容、来源、分类、关键词、向量特征、发布时间和互动计数，作为推荐候选内容池。
  \item 推荐Feed生成。系统需要根据用户身份、请求参数、语言、地区、历史曝光和行为特征生成个性化Feed。
  \item 行为埋点。系统需要记录曝光、点击、停留、分享等事件，并将其写入云端数据库，用于后续兴趣更新和算法评估。
  \item 实验配置。系统需要支持不同推荐策略的灰度实验，通过实验配置控制流量比例、分桶方式和评价指标。
  \item 接口联调。系统需要通过Apifox完成关键接口的请求配置、参数检查和测试记录。
\end{enumerate}

\subsection{非功能需求}

除功能需求外，系统还应满足以下非功能需求：

\begin{enumerate}
  \item 可扩展性。推荐召回源、过滤器和评分器应支持独立扩展，避免算法调整影响全部后端逻辑。
  \item 低延迟。Feed接口属于高频接口，应通过缓存、超时控制和候选数量限制降低响应时间。
  \item 可解释性。每条推荐结果应尽量能够追踪其召回来源、主要评分因素和实验分桶。
  \item 安全性。涉及用户资料、密码字段和行为数据的表必须加强访问控制，避免匿名访问造成数据泄露。
  \item 可测试性。接口路径、请求参数和响应结构应能在Apifox中清晰表达，便于课程验收和后续自动化测试。
\end{enumerate}

\section{可行性分析}

\subsection{技术可行性}

Spring Boot提供成熟的Web开发、依赖注入、数据访问和安全控制能力，能够满足后端服务搭建需求。Supabase PostgreSQL已经具备项目所需的核心表结构，能够保存用户、内容、行为和实验数据。Redis适合处理时间线缓存和热点Feed缓存，能够降低数据库压力。Apifox可以完成接口联调和截图取证。因此，从技术栈角度看，本项目具备实现可行性。

\subsection{经济可行性}

项目采用开源框架和云端免费层服务完成课程设计。Spring Boot、React、Redis客户端和大部分推荐算法逻辑均可免费使用。Supabase项目当前处于免费计划，能够满足课程阶段的数据规模。Apifox桌面端可用于本地联调，不需要额外购买商业服务。

\subsection{操作可行性}

系统采用前后端分离方式，用户通过前端页面访问推荐内容，开发者通过Apifox测试接口，管理员可通过Supabase控制台查看数据库。对于课程展示而言，操作路径清晰，既可以展示数据库表结构，也可以展示接口请求配置和推荐算法流程。

\section{项目系统架构}

系统总体架构如图\ref{fig:system_architecture_springboot}所示。整体分为客户端层、业务服务层、异步投递层、图计算层、数据与缓存层、推荐算法层六部分。客户端层包括React前端和Apifox接口联调工具；业务服务层以Spring Boot为核心，负责REST接口、认证与权限、推荐Feed服务、用户与群组服务、新闻内容服务；异步投递层由Go服务负责Redis Stream消费、消息扇出、outbox状态维护和同步唤醒；图计算层由C++服务负责社交图谱快照、邻居查询、共同关系查询和多跳遍历；数据与缓存层包括Redis、Supabase PostgreSQL、MongoDB和对象存储；推荐算法层负责多路召回、过滤规则、加权评分、TopK排序和埋点实验。

\begin{figure}[htbp]
  \centering
  \includegraphics[width=0.95\textwidth]{system-architecture-springboot.png}
  \caption{系统总体架构图}
  \label{fig:system_architecture_springboot}
\end{figure}

后端入口采用REST API。推荐Feed接口可以设计为：

\begin{lstlisting}[language=Java,caption={推荐Feed接口设计示例}]
@RestController
@RequestMapping("/api/space")
public class SpaceFeedController {
    private final RecommendationService recommendationService;

    @GetMapping("/feed")
    public FeedResponse feed(
            @AuthenticationPrincipal LoginUser user,
            @RequestParam(defaultValue = "20") int limit,
            @RequestParam(required = false) String cursor,
            @RequestParam(defaultValue = "false") boolean inNetworkOnly) {
        FeedRequest request = FeedRequest.of(user.getId(), limit, cursor, inNetworkOnly);
        return recommendationService.recommend(request);
    }
}
\end{lstlisting}

该接口将参数解析与业务处理分离。Controller只负责接收请求和校验参数，具体推荐流程交给\texttt{RecommendationService}完成。

\section{项目设计}

\subsection{功能模块设计}

系统按照职责划分为以下模块：

\begin{longtable}{p{0.22\textwidth}p{0.7\textwidth}}
\caption{系统功能模块设计}\\
\toprule
模块 & 主要职责 \\
\midrule
\endfirsthead
\toprule
模块 & 主要职责 \\
\midrule
\endhead
用户模块 & 管理用户资料、登录态、地区、语言、在线状态，为推荐服务提供用户基础画像。\\
关系模块 & 管理联系人、关注关系、群组成员关系，为社交召回和作者亲和度计算提供依据。\\
内容模块 & 管理新闻内容、分类、关键词、向量、发布时间、来源和互动计数。\\
推荐模块 & 编排特征补全、多路召回、候选补全、过滤、评分、TopK选择和曝光记录。\\
Go投递模块 & 消费Redis Stream中的平台事件和消息事件，按业务键并行投递，维护outbox聚合状态，批量ACK并处理失败恢复。\\
C++图计算模块 & 加载社交图谱快照，提供直接邻居、共同关系、ranked邻居和多跳遍历查询，为推荐召回提供低延迟图特征。\\
实验模块 & 维护实验配置、流量比例、实验分桶和评价指标，支持推荐策略灰度发布。\\
测试模块 & 通过Apifox维护接口请求、参数、返回结构和联调截图。\\
\bottomrule
\end{longtable}

\subsection{Go消息投递链路设计}

即时通信系统中的消息投递具有明显的分布式特征：同一会话内的消息需要保持顺序，不同会话之间又需要尽可能并行；投递过程中可能出现Redis读取失败、MongoDB写入超时、接收者状态变化和消费者重启等情况。因此，本项目将投递链路从业务接口中拆分出来，由Go服务独立负责。

Go投递链路的数据流如下：

\begin{enumerate}
  \item Spring Boot业务服务或平台事件总线将消息事件写入Redis Stream。
  \item Go消费者按consumer group读取消息，先解析事件类型和业务分区键。
  \item 调度器根据\texttt{chat\_id}、\texttt{user\_id}、\texttt{outbox\_id}或\texttt{message\_id}选择worker lane。
  \item worker在lane内串行处理同一业务键，执行接收者投影、update id预留、outbox chunk更新和wake事件发布。
  \item 成功消息按stream聚合后批量\texttt{XACK}，失败消息进入重试或死信处理。
\end{enumerate}

该设计的关键是“局部有序、全局并行”。设消息集合为$\mathcal{M}$，业务分区函数为$k(m)$，worker lane数量为$W$，则消息$m$被映射到：

\[
lane(m)=hash(k(m)) \bmod W
\]

对于任意两条消息$m_i,m_j$，若$k(m_i)=k(m_j)$，则二者进入同一lane并按读取顺序处理；若$k(m_i)\ne k(m_j)$，则可以进入不同lane并行处理。这样既不会破坏同一聊天或同一outbox内部的顺序，也避免所有消息被单线程消费限制吞吐。

在写入侧，Go服务使用分块预留降低MongoDB远程写放大。传统做法可能为每个接收者单独申请一个update id，若一个群消息有$n$个接收者，则远程写次数接近$O(n)$。分块预留将多个接收者合并为block，一次远程操作申请连续区间，再在本地分配给接收者，使远程预留次数近似降为$O(\lceil n/B \rceil)$，其中$B$为block大小。outbox状态也不再每次读回所有chunk重算，而是通过增量计数维护\texttt{queuedChunkCount}、\texttt{completedChunkCount}和\texttt{failedChunkCount}，再保留低频reconcile路径用于修复异常状态。

为了支持运维和课程验收，Go服务在summary中记录worker lane数量、ack batch size、batch ack count、reservation mode、wake publish mode、outbox aggregate mode、Mongo失败分类和主投递成功率等字段。通过这些字段可以判断投递链路是否处于高并发模式、是否发生批量ACK、是否出现大量可重试失败，以及是否需要调整worker数量或批量大小。

\subsection{C++图计算内核设计}

推荐系统中的社交召回依赖图结构。用户、联系人、群组、作者和互动关系可以抽象为图$G=(V,E)$，其中$V$表示用户或实体节点，$E$表示关注、联系人、共同群组、互动等边。若在每次推荐请求中动态查询数据库计算共同关系，延迟会受到网络往返和数据库负载影响。C++图计算内核通过内存快照方式，将高频图查询从业务数据库路径中分离出来。

图快照构建分为三个阶段。第一阶段从后端快照接口按页读取边数据，避免一次性加载全部边造成内存峰值。第二阶段将字符串形式的\texttt{source\_user\_id}、\texttt{target\_user\_id}和边类型intern为\texttt{uint32}编号。第三阶段按源点聚合并构建CSR索引：

\[
neighbors(src)=N[offset[src]\,..\,offset[src+1])
\]

其中$offset$数组记录每个源点邻居区间的起止位置，$N$数组连续保存邻居编号、边权和时间信息。该结构适合直接邻居读取和ranked top-K读取，也便于在overlap查询中顺序扫描两个邻居列表。

共同关系查询可表示为两个邻居集合的交集：

\[
overlap(u,v)=N(u)\cap N(v)
\]

如果先物化全部交集再排序，在高出度节点上会消耗大量内存。C++服务采用streaming top-K overlap：扫描邻居列表时只维护得分最高的$K$个共同节点，并记录scanned、visited和budget exhausted等指标。多跳遍历同样采用预算控制，避免在社交图高度连接时无限扩展。对于排序中需要的新鲜度和时间衰减，服务在请求入口只采样一次\texttt{now\_ms}，并在ranking、traversal和weight计算中复用，减少热循环系统调用。

图服务对外提供HTTP/JSON接口，主要包括健康检查、邻居查询、共同关系查询和诊断summary。健康接口返回快照是否加载成功，summary返回快照版本、快照表示、节点数、边数、构建阶段耗时和快照年龄。推荐服务可以根据这些字段判断是否启用图召回、是否降级到热门召回，以及是否需要等待下一轮快照刷新。

\subsection{推荐算法流程设计}

推荐算法流程如图\ref{fig:recommendation_pipeline}所示。推荐请求进入系统后，首先构造请求上下文，包括用户编号、分页游标、语言、地区、客户端类型和已曝光内容。随后系统对用户画像进行补全，包括关注关系、兴趣向量、历史行为、实验桶和近期请求状态。

\begin{figure}[htbp]
  \centering
  \includegraphics[width=0.98\textwidth]{recommendation-pipeline.png}
  \caption{个性化推荐算法流程图}
  \label{fig:recommendation_pipeline}
\end{figure}

\subsubsection{用户-内容交互建模}

推荐系统的本质是对用户$u \in \mathcal{U}$与内容$c \in \mathcal{C}$之间的潜在偏好进行估计。设用户在时刻$t$对内容$c$产生的行为为$y_{u,c,t}$，行为类型集合为

\[
\mathcal{A}=\{\text{impression}, \text{click}, \text{dwell}, \text{share}, \text{negative}\}
\]

其中曝光代表系统曾把内容展示给用户，点击、停留、分享属于正向反馈，不感兴趣、屏蔽、举报属于负向反馈。由于真实系统中绝大多数内容不会被用户显式评分，本项目采用隐式反馈建模思想，将行为强度映射为偏好置信度。对用户$u$和内容$c$，定义隐式偏好值：

\[
r_{u,c}=\sum_{a \in \mathcal{A}} \alpha_a \cdot n_{u,c}^{(a)}
\]

其中$n_{u,c}^{(a)}$表示行为$a$出现的次数，$\alpha_a$表示该行为的权重。点击权重较低，表示“有兴趣查看”；长停留、分享、回复权重较高，表示“深度兴趣”；负反馈权重为负，表示系统应该降低相似内容的排序。该思想与隐式反馈推荐和排序学习中的偏好强度建模一致\cite{hu2008collaborative,rendle2009bpr}。

为了同时表达用户的短期兴趣和长期兴趣，系统为每个用户维护两个向量：

\[
\mathbf{p}_u = \lambda \mathbf{p}_u^{short} + (1-\lambda)\mathbf{p}_u^{long}
\]

其中$\mathbf{p}_u^{short}$由最近点击、停留、分享的内容向量加权得到，更强调当前会话和近期兴趣；$\mathbf{p}_u^{long}$由较长时间窗口内的稳定行为得到，更强调长期偏好。$\lambda$是时间敏感系数，当用户近期行为密集时增大，当近期行为稀疏时减小。内容$c$的表示向量记为$\mathbf{q}_c$，可以由标题、摘要、关键词、分类和已有embedding字段综合得到。用户与内容的语义匹配分数定义为：

\[
s_{match}(u,c)=\cos(\mathbf{p}_u,\mathbf{q}_c)=
\frac{\mathbf{p}_u^\top \mathbf{q}_c}{\|\mathbf{p}_u\|_2\|\mathbf{q}_c\|_2}
\]

该得分用于向量召回和排序特征，能够弥补单纯社交关系召回无法发现新兴趣内容的问题。

\subsubsection{工程实现中的用户推荐特征清单}

为了便于复检，本项目将推荐算法使用的用户特征分为“持久化画像向量、实时行为信号、长期行为日志、社交图谱边、请求级上下文、候选评分特征”六类。前端只负责采集用户真实动作，后端再将动作转换为\texttt{user\_actions}、\texttt{user\_signals}、\texttt{real\_graph\_edges}和\texttt{user\_feature\_vectors}等推荐特征，不允许前端直接伪造兴趣向量或模型分数。

\begingroup
\hbadness=10000
\small
\begin{longtable}{L{0.17\textwidth}L{0.33\textwidth}L{0.39\textwidth}}
\toprule
特征层 & 主要字段 & 用途说明\\
\midrule
\endhead
持久化用户向量 & \texttt{interestedInClusters}、\texttt{knownForCluster}、\texttt{knownForScore}、\texttt{producerEmbedding} & 表示用户消费兴趣、生产者社区归属和作者侧稀疏向量，用于SimClusters式兴趣召回和作者匹配。\\
稠密Embedding & \texttt{twoTowerEmbedding}、\texttt{phoenixEmbedding}、\texttt{twhinEmbedding}、\texttt{embeddingDim} & 表示双塔、Phoenix和图嵌入空间中的用户向量，用于ANN召回、语义匹配和排序特征。\\
向量契约元数据 & \texttt{version}、\texttt{modelVersion}、\texttt{artifactVersion}、\texttt{modelProfile}、\texttt{embeddingContract}、\texttt{computedAt}、\texttt{expiresAt}、\texttt{qualityScore} & 记录向量由哪个模型和artifact生成、维度是否兼容、何时计算和是否过期，避免不同维度或不同版本向量混算。\\
实时用户信号 & \texttt{signalType}、\texttt{targetType}、\texttt{targetId}、\texttt{targetAuthorId}、\texttt{productSurface}、\texttt{metadata}、\texttt{timestamp}、\texttt{expiresAt} & 保存短期兴趣与负反馈，例如点击、停留、点赞、转发、引用、关注、不感兴趣、屏蔽、静音和举报，供实时特征聚合使用。\\
长期行为日志 & \texttt{action}、\texttt{targetPostId}、\texttt{targetAuthorId}、\texttt{requestId}、\texttt{rank}、\texttt{score}、\texttt{weightedScore}、\texttt{recallSource}、\texttt{selectionPool}、\texttt{targetKeywords}、\texttt{targetUrl} & 保存可回放、可审计的行为事实，用于训练样本、行为序列、推荐解释和线上问题追踪。\\
RealGraph社交边 & \texttt{followCount}、\texttt{likeCount}、\texttt{replyCount}、\texttt{retweetCount}、\texttt{quoteCount}、\texttt{tweetClickCount}、\texttt{dwellTimeMs}、\texttt{muteCount}、\texttt{blockCount}、\texttt{reportCount}、\texttt{decayedSum}、\texttt{interactionProbability} & 将用户与作者之间的互动累积为图边特征，正向互动增强作者亲和，静音、屏蔽、举报和取关降低或阻断关系分数。\\
请求级用户上下文 & \texttt{followedUserIds}、\texttt{blockedUserIds}、\texttt{mutedUserIds}、\texttt{mutedKeywords}、\texttt{seenPostIds}、\texttt{followerIds}、\texttt{userStateContext}、\texttt{userSignalFeatures}、\texttt{interestedTopics} & 在每次Feed请求开始时由Hydrator补齐，用于关注召回、屏蔽过滤、静音过滤、已看过滤、冷启动/稀疏/重度用户分层和兴趣话题匹配。\\
候选评分特征 & \texttt{likeScore}、\texttt{replyScore}、\texttt{repostScore}、\texttt{quoteScore}、\texttt{clickScore}、\texttt{profileClickScore}、\texttt{dwellScore}、\texttt{followAuthorScore}、\texttt{notInterestedScore}、\texttt{blockAuthorScore}、\texttt{muteAuthorScore}、\texttt{reportScore} & 对每条候选内容生成动作概率和负反馈概率。加权评分器对正向行为加分，对不感兴趣、屏蔽作者、静音作者和举报进行扣分。\\
\bottomrule
\end{longtable}
\endgroup

其中，\texttt{UserSignal}中的正权重信号包括点赞、转发、回复、引用、关注、分享、帖子点击、主页点击、外链点击、话题点击、搜索查询、曝光和停留；\texttt{QUOTE}按$+2.5$进入实时信号，并汇总为\texttt{UserSignalFeatures.quoteCount}和RealGraph的\texttt{quoteCount}。负权重信号包括取消点赞、取消转发、取消关注、屏蔽、静音、举报、取消书签和通知关闭，对应权重分别为$-0.5$、$-1.0$、$-3.0$、$-10.0$、$-5.0$、$-8.0$、$-0.5$和$-0.1$。\texttt{DISMISS\_POST}和\texttt{HIDE\_POST}虽然当前不在通用权重表中直接给出固定分值，但会写入\texttt{user\_actions}和\texttt{user\_signals}，并在候选负反馈特征、训练样本和Feed过滤/降权逻辑中使用。

\begingroup
\hbadness=10000
\scriptsize
\begin{longtable}{L{0.13\textwidth}L{0.20\textwidth}L{0.27\textwidth}L{0.26\textwidth}}
\toprule
正向/兴趣特征 & 前端入口 & 落库/接口 & 推荐侧影响\\
\midrule
\endhead
曝光与停留 & 帖子进入视口和可见停留计时 & \texttt{impression/dwell}埋点桥接到\path{user_actions}和\path{user_signals} & 形成served-history、行为序列和隐式兴趣强度。\\
帖子点击与主页点击 & 点击帖子卡片、头像或用户名 & \texttt{click/profile_click}埋点写入\texttt{ActionType.CLICK/PROFILE_CLICK}和对应\texttt{SignalType} & 增强内容兴趣和作者亲和，供Hydrator读取近期行为。\\
点赞、回复、转发 & 帖子操作栏按钮和评论抽屉 & \path{/like}、\path{/comments}、\path{/repost}写入推荐事件 & 增强显式兴趣、训练标签和RealGraph作者边。\\
引用动态 & 帖子更多菜单中的“引用动态”按钮，输入引用内容后发布 & \path{POST /api/space/posts}携带\texttt{quotePostId/quoteContent}，写入\texttt{ActionType.QUOTE}和\texttt{SignalType.QUOTE} & 进入\texttt{quoteCount}、\texttt{quoteScore}、Phoenix训练标签和RealGraph社交边。\\
分享与外链点击 & 分享按钮、Markdown外链 & 分享入口和\texttt{open_link}埋点写入\texttt{SignalType.SHARE/OPEN_LINK}，外链保留\texttt{targetUrl} & 作为轻量兴趣和内容传播信号。\\
搜索与话题点击 & 探索页搜索框、热门趋势和话题chip & \texttt{search_query/hashtag_click}桥接为\texttt{TargetType.SEARCH_QUERY/TOPIC} & 更新兴趣关键词和探索面意图。\\
关注 & 推荐关注卡片或个人主页关注按钮 & \path{POST /api/space/users/:id/follow}和analytics对账写入\texttt{SignalType.FOLLOW} & 进入关注召回、作者亲和和RealGraph。\\
\bottomrule
\end{longtable}
\endgroup

\begingroup
\hbadness=10000
\scriptsize
\begin{longtable}{L{0.13\textwidth}L{0.19\textwidth}L{0.25\textwidth}L{0.29\textwidth}}
\toprule
负向特征 & 前端入口 & 落库/接口 & 推荐侧影响\\
\midrule
\endhead
不感兴趣 & 帖子更多菜单中的“不感兴趣”按钮 & \path{/api/space/posts/:id/not-interested}，写入\path{ActionType.DISMISS}和\path{SignalType.DISMISS_POST} & 作为帖子级强负反馈，进入\path{notInterestedScore/dismissScore}和后续训练样本，降低相似内容排序。\\
隐藏帖子 & 帖子更多菜单中的“隐藏此帖”按钮 & \path{/api/space/posts/:id/hide}，写入\path{ActionType.HIDE}和\path{SignalType.HIDE_POST} & 当前请求中移除该帖子，并作为负反馈事实用于后续过滤、行为审计和训练。\\
举报 & 帖子更多菜单中的“举报”按钮，二级原因包含垃圾内容、骚扰、虚假信息、暴力仇恨和其他原因 & \path{/api/space/posts/:id/report}，写入\texttt{ActionType.REPORT}和\texttt{SignalType.REPORT} & 在\texttt{WeightedScorer}中按\texttt{REPORT\_WEIGHT=-8.0}扣分，同时进入RealGraph的\texttt{reportCount}。\\
屏蔽作者 & 帖子更多菜单中的“屏蔽 @作者”按钮 & \path{/api/space/users/:id/block}，写入联系人\texttt{blocked}状态和\texttt{SignalType.BLOCK} & 进入\texttt{blockedUserIds}硬过滤、\path{blockAuthorScore/blockScore}扣分和RealGraph的\texttt{blockCount}。\\
静音作者 & 帖子更多菜单中的“静音 @作者”按钮 & \path{/api/space/users/:id/mute}，写入用户设置和\texttt{SignalType.MUTE} & 进入\texttt{mutedUserIds}过滤、\texttt{muteAuthorScore}扣分和RealGraph的\texttt{muteCount}。\\
取消关注 & 推荐关注卡片或个人主页中的关注状态按钮，已关注时显示为“已关注”并可再次点击取消 & \texttt{DELETE} \path{/api/space/users/:id/follow}，写入\texttt{SignalType.UNFOLLOW} & 将作者关系从关注网络移除，并在RealGraph中减少\texttt{followCount}。\\
取消点赞 & 帖子点赞按钮的已点赞状态再次点击 & 前端埋点\texttt{unlike}映射到\path{SignalType.UNFAVORITE}，业务接口为\texttt{DELETE} \path{/api/space/posts/:id/like} & 作为轻度负向修正，降低用户对该内容/相似内容的偏好置信度。\\
取消转发 & 帖子转发按钮在已转发状态下显示为“取消转发”语义按钮 & \texttt{DELETE} \path{/api/space/posts/:id/repost}，写入\path{SignalType.UNRETWEET} & 与\texttt{UNRETWEET=-1.0}负权重对齐，作为撤销转发的轻度负反馈进入实时用户信号。\\
取消书签、通知关闭 & \texttt{UserSignal}权重表预留 & 当前Space Feed未提供书签和通知关闭入口 & 属于其他业务面或后续扩展字段，不应在Space演示中宣称已经由前端按钮采集。\\
\bottomrule
\end{longtable}
\endgroup

2026年6月11日通过Chrome DevTools MCP在生产前端\texttt{/space}页面进行只读复检：打开非本人作者\texttt{demo\_growth\_author\_06}的帖子更多菜单后，页面DOM中实际存在“不感兴趣”“隐藏此帖”“举报”“屏蔽 @demo\_growth\_author\_06”“静音 @demo\_growth\_author\_06”五个按钮；进入举报二级菜单后，实际存在“垃圾内容/广告”“骚扰/霸凌”“虚假信息”“暴力/仇恨”“其他原因”五个原因按钮。本次复检没有点击最终提交类按钮，网络请求中未出现\texttt{/not-interested}、\texttt{/hide}、\texttt{/report}、\texttt{/block}或\texttt{/mute}写接口，因此不会污染线上用户特征。取消转发闭环在代码和组件测试中补齐：当帖子处于已转发状态时，同一转发按钮的语义变为“取消转发”，并调用\texttt{DELETE /api/space/posts/:id/repost}写入\texttt{UNRETWEET}。引用闭环在本轮补齐：更多菜单新增“引用动态”入口，前端提交引用文本后调用\texttt{POST /api/space/posts}并携带\texttt{quotePostId}，后端写入\texttt{QUOTE}推荐事件。当前前端没有“拉踩”或\texttt{downvote/dislike}按钮，对应算法也未定义独立拉踩特征；演示时应将“不感兴趣、隐藏、举报、屏蔽、静音、取消转发、取消点赞、取消关注”作为负向减分入口。

\subsubsection{多路召回模型}

多路召回阶段的目标是在可控延迟内从海量内容集合$\mathcal{C}$中构造较小候选集$\mathcal{R}_u$。如果直接对所有内容排序，时间复杂度为$O(|\mathcal{C}|)$，在实际Feed场景中不可接受。因此系统将召回看作多个弱检索器的并集：

\[
\mathcal{R}_u = \bigcup_{m=1}^{M} \mathcal{R}_u^{(m)}
\]

其中$m$表示召回通道，$\mathcal{R}_u^{(m)}$表示第$m$路召回结果。每一路召回承担不同的信息增益：

\begin{enumerate}
  \item \textbf{关注召回}：从用户关注对象的时间线中读取近期内容，适合关注关系稳定的用户。
  \item \textbf{社交图谱召回}：根据共同联系人、共同群组和互动关系扩展候选作者。
  \item \textbf{新闻向量召回}：利用内容关键词和向量特征查找相似新闻。
  \item \textbf{作者兴趣召回}：根据用户长期兴趣向量匹配更可能感兴趣的作者。
  \item \textbf{热门召回}：选取近期点击、分享、停留表现较好的内容。
  \item \textbf{冷启动召回}：当用户行为较少时，使用地区、语言和全局热度生成候选。
\end{enumerate}

从概率角度看，每一路召回都近似估计一个条件概率$P(c|u,m)$。关注召回强调$P(c|follow(u))$，社交图谱召回强调$P(c|G_u)$，向量召回强调$P(c|\mathbf{p}_u)$，热门召回强调$P(c|global)$。最终候选集并不是简单平均，而是根据用户状态动态调整各路配额。设用户状态为$z_u \in \{\text{cold},\text{sparse},\text{warm},\text{heavy}\}$，第$m$路召回配额为：

\[
K_m(u)=K \cdot \frac{\pi_m(z_u)\cdot h_m}{\sum_{j=1}^{M}\pi_j(z_u)\cdot h_j}
\]

其中$K$表示总候选规模，$\pi_m(z_u)$表示不同用户状态下的召回偏好，$h_m$表示召回源健康度。对于冷启动用户，热门召回和地区语言召回权重更高；对于重度用户，关注召回、向量召回和作者亲和召回权重更高。该设计可以避免“新用户无结果”和“老用户结果过于热门化”两个常见问题。

\subsubsection{规则过滤与约束建模}

召回阶段追求覆盖率，过滤阶段追求安全性和有效性。设召回候选为$\mathcal{R}_u$，过滤后候选为$\mathcal{F}_u$，则：

\[
\mathcal{F}_u=\{c \in \mathcal{R}_u \mid g_k(u,c)=1,\; k=1,2,\dots,L\}
\]

其中$g_k$是第$k$个硬约束函数。系统中主要约束包括：去重约束、已曝光过滤、自身内容过滤、屏蔽用户过滤、敏感词过滤、内容时效过滤和重复转发过滤。硬约束的优先级高于模型得分，因为违反这些规则会直接损害用户体验或平台安全。

对已曝光内容，系统维护集合$\mathcal{E}_u$，如果$c \in \mathcal{E}_u$，则过滤该内容或显著降低得分。该机制解决推荐系统中的重复曝光问题。对负反馈内容，系统不仅过滤具体内容，还会将负反馈传播到相似作者、相似主题和相似关键词：

\[
penalty(u,c)=\beta_1 I(author(c)\in B_u)+\beta_2 sim(topic(c),T_u^{-})+\beta_3 sim(keyword(c),K_u^{-})
\]

其中$B_u$表示用户屏蔽作者集合，$T_u^{-}$表示负反馈主题集合，$K_u^{-}$表示负反馈关键词集合。该惩罚项进入排序阶段，避免系统反复推荐用户明确拒绝的内容。

\subsubsection{排序评分函数}

评分阶段采用可解释的多目标融合模型。设候选内容$c$对用户$u$的最终得分为$S(u,c)$，可表示为：

\[
S(u,c)=w_1E(u,c)+w_2A(u,c)+w_3Q(c)+w_4R(c)+w_5M(u,c)-w_6N(u,c)
\]

其中，$E(u,c)$表示用户行为兴趣得分，$A(u,c)$表示作者亲和度，$Q(c)$表示内容质量，$R(c)$表示新鲜度，$M(u,c)$表示语义匹配得分，$N(u,c)$表示负反馈惩罚项。权重$w_i$可以由实验配置控制，也可以在后续通过学习排序模型自动学习。

各部分含义如下：

\begin{enumerate}
  \item \textbf{行为兴趣得分}：由点击、停留、分享、回复等行为加权得到，表达用户对相似内容的历史偏好。
  \item \textbf{作者亲和度}：由关注关系、共同群组、历史互动、作者内容点击率等特征构成，衡量用户与作者之间的社交距离。
  \item \textbf{内容质量}：由内容完整度、来源可信度、历史互动率和负反馈率构成，避免低质量内容因短期热度被过度放大。
  \item \textbf{新鲜度}：使用时间衰减函数计算，保证Feed具备实时性。
  \item \textbf{语义匹配}：使用用户兴趣向量与内容向量余弦相似度计算，表达主题层面的匹配程度。
  \item \textbf{负反馈惩罚}：由屏蔽、不感兴趣、举报、关键词厌恶等信号组成，用于降低不符合用户偏好的内容。
\end{enumerate}

新鲜度采用指数衰减模型：

\[
R(c)=\exp\left(-\frac{t_{now}-t_c}{\tau}\right)
\]

其中$t_c$为内容发布时间，$\tau$为时间半衰参数。新闻类内容对新鲜度敏感，$\tau$应较小；知识类或长期有效内容对新鲜度不敏感，$\tau$可以增大。该函数可以避免旧内容长期占据高位，同时不会让刚发布但质量很低的内容无条件排在前面。

行为兴趣得分采用归一化加权形式：

\[
E(u,c)=\sigma\left(\sum_{a \in \mathcal{A}^{+}}\alpha_a f_a(u,c)-\sum_{b \in \mathcal{A}^{-}}\gamma_b f_b(u,c)\right)
\]

其中$\sigma(x)=1/(1+e^{-x})$为Sigmoid函数，$\mathcal{A}^{+}$为正向行为集合，$\mathcal{A}^{-}$为负向行为集合。使用Sigmoid的原因是将原始行为强度压缩到$[0,1]$区间，避免某个高频行为单独支配排序。

\subsubsection{排序校准与多样性重排}

如果直接按照$S(u,c)$取TopK，系统容易出现“信息茧房”和“作者集中”问题。为此，最终选择阶段采用带多样性约束的贪心重排。设已选择集合为$\mathcal{S}$，下一条内容$c^\ast$满足：

\[
c^\ast=\arg\max_{c \in \mathcal{F}_u \setminus \mathcal{S}}
\left[S(u,c)-\eta \max_{c' \in \mathcal{S}} sim(c,c')\right]
\]

其中$sim(c,c')$表示两个内容在作者、主题、来源或向量空间中的相似度，$\eta$为多样性惩罚系数。该公式与最大边际相关性（MMR）思想一致\cite{carbonell1998mmr}。当已选择列表中存在同主题内容时，相似候选会被惩罚，从而提升推荐列表的信息覆盖面。

同时，系统对不同召回来源进行校准。由于热门召回天然具有较高互动数，向量召回天然更贴近兴趣但可能互动较少，若不校准会导致某一路召回长期占优。系统引入来源校准项：

\[
S'(u,c)=S(u,c)+\delta_{source(c)}+\delta_{bucket(u)}
\]

其中$\delta_{source(c)}$用于修正不同召回源的系统性偏差，$\delta_{bucket(u)}$用于支持实验分桶。这样既可以做工程上的灰度实验，也可以在报告层面体现推荐算法的可观测和可调优能力。

\subsubsection{目标函数与评价指标}

从优化目标看，推荐系统不是单纯最大化点击率，而是多目标优化问题。系统希望同时提升点击、停留、分享和满意度，并降低负反馈。因此可定义总体效用函数：

\[
J=\mathbb{E}_{(u,c)\sim \mathcal{D}}
\left[\theta_1 click+\theta_2 dwell+\theta_3 share-\theta_4 negative-\theta_5 latency\right]
\]

其中$\mathcal{D}$为线上请求分布，$latency$表示推荐接口延迟。该目标函数体现了课程项目中的工程权衡：推荐结果要相关，但接口也要稳定、快速、可解释。离线评估可以使用Precision@K、Recall@K、NDCG@K和覆盖率；在线评估可以使用CTR、平均停留时长、分享率、负反馈率和P95延迟。实验表中的\texttt{experiments}表正是为后续A/B测试准备的。

综合来看，本项目推荐算法并非只按热度排序，而是形成了“特征建模--多路召回--规则过滤--多目标排序--多样性重排--实验评估”的完整闭环。该闭环与工业推荐系统和经典深度推荐模型的设计思想一致\cite{covington,cheng,he2017ncf}，同时保留了课程项目需要的可解释性与工程可实现性。

\subsection{数据库设计}

通过Supabase MCP读取可知，本项目云端数据库项目为\texttt{telegram}，核心表结构如图\ref{fig:supabase_er_diagram}所示。数据库主要分为基础社交模块、推荐内容模块和实验平台模块。

\begin{figure}[htbp]
  \centering
  \includegraphics[width=0.95\textwidth]{supabase-er-diagram.png}
  \caption{Supabase数据库实体关系图}
  \label{fig:supabase_er_diagram}
\end{figure}

\begin{longtable}{L{0.25\textwidth}L{0.18\textwidth}L{0.48\textwidth}}
\caption{核心数据表说明}\\
\toprule
表名 & 当前记录数 & 用途说明 \\
\midrule
\endfirsthead
\toprule
表名 & 当前记录数 & 用途说明 \\
\midrule
\endhead
\texttt{users} & 642 & 保存用户编号、用户名、邮箱、头像、地区、语言和在线状态。\\
\texttt{groups} & 6 & 保存群组名称、描述、群主、类型、成员数量和最大成员数。\\
\texttt{group\_members} & 764 & 保存用户与群组的成员关系、角色、状态和加入时间。\\
\texttt{contacts} & 195 & 保存联系人关系、联系人状态和备注名。\\
\texttt{news\_articles} & 838 & 保存新闻标题、摘要、来源、分类、关键词、向量、互动计数和有效状态。\\
\texttt{news\_user\_events} & 2386 & 保存曝光、点击、停留、分享等用户行为事件。\\
\texttt{news\_user\_vectors} & 6 & 保存用户短期兴趣向量和长期兴趣向量。\\
\texttt{experiments} & 2 & 保存推荐实验名称、状态、流量比例、分桶方式和评价指标。\\
\bottomrule
\end{longtable}

\subsection{接口设计}

系统核心接口设计如下：

\begin{longtable}{L{0.13\textwidth}L{0.34\textwidth}L{0.41\textwidth}}
\caption{核心接口设计}\\
\toprule
方法 & 路径 & 说明 \\
\midrule
\endfirsthead
\toprule
方法 & 路径 & 说明 \\
\midrule
\endhead
GET & \path{/api/space/feed} & 获取个性化推荐Feed，支持分页、数量限制和是否只看关注内容。\\
POST & \path{/api/space/feed} & 携带已曝光内容、请求编号和客户端上下文获取Feed。\\
GET & \path{/api/space/trends} & 获取趋势主题和热门内容。\\
GET & \path{/api/space/recommend/users} & 推荐可能感兴趣的用户或作者。\\
POST & \path{/api/events/news} & 上报曝光、点击、停留、分享等行为事件。\\
\bottomrule
\end{longtable}

\subsection{安全设计}

安全设计主要包括接口鉴权、参数校验、数据库访问控制和敏感字段保护。Spring Boot后端使用JWT识别当前用户，推荐接口只允许登录用户访问。请求参数如\texttt{limit}需要设置上限，避免一次请求拉取过多候选内容。涉及密码、用户行为和兴趣向量的数据不应直接暴露给前端。

需要特别说明的是，Supabase安全检查发现当前公开模式下多张表尚未启用RLS，包括\texttt{users}、\texttt{groups}、\texttt{contacts}、\texttt{news\_articles}、\texttt{news\_user\_events}等。课程报告中将其作为部署安全风险记录，实际上线前应启用RLS并补充按用户、角色和服务端密钥区分的访问策略。

\section{本章小结}

本章完成了项目的需求分析、可行性分析、架构设计、模块设计、推荐算法流程设计、数据库设计和接口设计。系统以Spring Boot作为后端核心，以Supabase PostgreSQL作为云端数据库，以Redis提升缓存性能，以Apifox支撑接口联调。推荐算法采用多路召回、规则过滤、加权评分和TopK选择的分层流程，能够满足课程大作业对分布式架构和推荐系统思想的要求。
